\section{Problem formulation and related work}\label{sec:prelim}

\subsection{Problem formulation}\label{sec:problem}

We write $\N=\{0,1,2,\ldots\}$ and use $\Z$ for the integers. An undirected temporal graph is $G=(V,\cE,\tau)$, with edge labels $\tau:\cE\to\N$, $n=\abs{V}$ and $m=\abs{\cE}$. Parallel edges are distinct. Loops are excluded. At time $t$, $G_t=(V,E_t)$, where $E_t=\{e\in\cE:\tau(e)=t\}$, is the \emph{snapshot}. A snapshot component contains at least one edge.

A \emph{non-strict temporal walk} is $W=(v_0,e_1,v_1,\ldots,e_\ell,v_\ell)$, where $e_i$ joins $v_{i-1}$ to $v_i$ and $\tau(e_1)\le\cdots\le\tau(e_\ell)$. Its start and end are $v_0,v_\ell$. An empty walk has $\ell=0$. Exact edge cover traverses each edge once across all walks.

The decision problems prescribe a vector $a\in\N^V$ of start counts, with $k=\sum_v a_v$.

\begin{definition}\label{def:problems}
\begin{itemize}
\item \FREE{}: decide whether there exist $k$ non-strict temporal walks that collectively traverse every edge of $\cE$ exactly once, with exactly $a_v$ walks starting at each vertex $v$.
\item \FIX{}: given an additional vector $z\in\N^V$, decide whether there exist walks satisfying \FREE{} with exactly $z_v$ walks ending at each vertex $v$.
\end{itemize}
\end{definition}

In \FIX{}, $a$ and $z$ prescribe endpoint counts without pairing starts to ends. Walks are non-strict unless stated otherwise.

Running times count unit-cost arithmetic operations on machine-word integers, including $a_v,z_v$. We return a \emph{compressed} solution (non-empty walks and each vertex's empty-walk count), since binary multiplicities can require exponentially many listed empty walks.

\subsection{Related work}\label{sec:related}

\paragraph{Exact edge cover in temporal graphs.}
Deligkas et al.~\cite{deligkas2025} study temporal paths, trails and walks (theorem and page numbers refer to the AAAI version, while the extended version is cited as arXiv v1). For strict walks with fixed terminals, Corollary~6.4 handles proper temporal graphs (each snapshot is a matching), and Theorem~6.5 gives an $O(2^k k\,t_{\max})$ algorithm when $|E_t|\in\{0,k\}$ ($t_{\max}$ is the largest time label). Directed cases are polynomial (Theorems~6.1 and~5.2). Their Theorem~6.3 mentions start terminals, while its proof and Figure~3 of~\cite{deligkas2025} use end terminals. The extended version~\cite{deligkas2024arxiv} states both. Section~\ref{sec:counterexample} analyses in detail only the end-terminal construction shown in that proof and in Figure~3 of~\cite{deligkas2025}. The variant with start terminals only is not addressed here.

\paragraph{Temporal networks and covering problems.}
Time-copy representations appear in~\cite{akrida2019}. Marino and Silva study a single walk that traverses every edge once in a temporal graph~\cite{marino2023}. For $k=1$ with a prescribed start and one label per edge, non-strict feasibility reduces, by Theorem~\ref{thm:char}, to a snapshot-by-snapshot Euler trail check, so we do not compare complexity results across the two models. Other related problems explore all edges~\cite{bumpus2023} or cover vertices of a temporal DAG~\cite{chakraborty2024}.

\paragraph{Discrete convexity and orientations.}
We use M-convex sets~\cite{murota1996,murota2003}, submodular flows~\cite{edmonds1977,franktardos1988,murota2003,fujishige2005}, degree-constrained orientations~\cite{hakimi1965,frank1980,frank2011} and augmenting paths~\cite{cunningham1985,fujishige2005}. Our additional constraint removes an even component's dead state. Related structures include strongly connected orientations~\cite{abdi2025}.
